\section{Conclusion}


For a long time, the dominant paradigm for building models in AI has centered around multiplication of tensor representations, and in the deep learning era this has given rise to highly modular (layer-wise) designs that allow for fast development and wide exploration. However, these design paradigms require extensive domain expertise, and even experts face substantial challenges when it comes to combining different pretrained components into larger systems.

The promise of in-context learning is that we can build complex systems from pretrained components using only natural language as the medium for giving systems instructions and, as we argue for, allowing components to communicate with each other. In this new paradigm, the building blocks are pretrained models and the core operations are natural language instructions and operations on natural language texts. If we can realize this potential, then we can broaden participation in AI system development, rapidly prototype systems for new domains, and maximize the value of specialized pretrained components.

In the current paper, we introduced the \demonstrate--\search--\predict (\DSP) framework for retrieval augmented in-context learning. \DSP{} consists of a number of simple, composable functions for implementing in-context learning systems as deliberate \textit{programs}---instead of end-task prompts---for solving knowledge intensive tasks. We implemented \DSP{} as a Python library and used it to write programs for Open-SQuAD, HotPotQA, and QReCC. These programs deliver substantial gains over previous in-context learning approaches. However, beyond any particular performance number, we argue that the central contribution of \DSP{} is in helping to reveal a very large space of conceptual possibilities for in-context learning in general.